\subsection{Definition}

(1) We say $f$ has \textbf{regular point} at $z_0$ if the singular part of the Laurent series at $z_0$ vanishes. \\
\textbf{Example :}
\begin{align*}
    f(z) = \sum_{u=0}^\infty \frac{z^u}{u!} = e^z, \text{ 0 regular}
\end{align*}
\textbf{Remark: } if $f$ is holomorphic at $z_0$ then $z_0$ is a regular point. Then converte is also true: if $z_0$ is a regular point pf $f$, then $\lim_{z \rightarrow z_0} f(z)$ exists and if we extend $f$ to $z_0$ this way, $f$ is holomorphic at $z_0$. In short $z_0$ regular point of $f$ if and only if $\lim_{z \rightarrow z_0} f(z)$ exists. \\ \\
(2) we say $f$ has a \textbf{pole of order $m$} at $z_0$ if the coefficient of the Laurent series of $f$ at $z_0$ satisfy $c_n = 0$ for $u < -m$, ie
\begin{align*}
    f(z) = \frac{c_{-m}}{(z-z_o)^m} + \frac{c_{-m+1}}{(z-z_0)^{m-1}} + \cdots, c_{-m} \neq 0 \\
    "f \text{ blows up with } \frac{1}{(z-z_0)^m } \text{ at } z_0"
\end{align*}
\textbf{Example}
\begin{itemize}
\item $f(z) = \frac{1}{z}$ has a pole of order $1$ at $z_0 = 0$
\item $f(z) = \frac{1}{(z + 2)^2} + \frac{1}{z}$ has a pole of order $2$ at $z_0 = -2$ and a pole of order $1$ at $z_0 = 0$
\end{itemize}
If $f$ has a pole of order $n$ at $z_0$, 
\begin{align*}
    f(z) = \frac{c_{-m}}{(z-z_0)^m} + \frac{c_{-m+1}}{(z-z_0)^{m-1}} + \cdots
\end{align*}
\textbf{Remarks} $g(z) = (z-z_0)^mf(z)$ has a regular point at $z_0$ if $f$ has a pole of order $m$ at $z_0$. \\
In fact $f$ has a pole of order $m$ at $z_0$ if and only if $\lim_{z \rightarrow z_0}(z-z_0)^mf(z)$ exist and is not zero. In particular if $f$ has a pole of order $z_0$ we can write $f(z) = \frac{g(z)}{(z-z_0)^m}$ where $g$ holomorphic at $z_0$ and $g(z_0) \neq 0.$ (opposite is also true) \\ \\
(3) we say $f$ has an essential singularity at $z_0$ if its Laurent series has infinitely many coefficient $c_n \neq 0$ with $n \leq -1$ \\
\textbf{Example :}
\begin{align*}
    f(z) = e^{1/z} \text{ at } z_0 = 0.
\end{align*}
\begin{align*}
    f(z) = \sum_{u=0}^\infty \frac{(\frac{1}{z})^u}{u!} = \sum_{u = - \infty}^0\frac{1}{(-u)!}z^u
\end{align*}
\textbf{Example :}
\begin{align*}
    f(z) = ze^{1/z} \text{ at } z_0 = 0.
\end{align*}
\begin{align*}
    f(z) = z  \sum_{u = - \infty}^0\frac{1}{(-u)!}z^u = z+1 + \frac{1}{2!z} + \frac{1}{3!z} + \cdots 
\end{align*}
So $0$ essentially singularity of $f$. \\
Now let's discuss quotients of Taylor series in this context. Suppose $p,q: D \setminus\{z_0\} \rightarrow \mathbb{C}$ holo with $z_0$ being regular for both. Write ther Laurent series as 
\begin{align*}
    p(z) = a_0 + a_1(z-z_0) + a_2(z-z_0)^2 + \cdots \\
    q(z) = b_0 + b_1(z-z_0) + b_2(z-z_0)^2 + \cdots 
\end{align*}
Goal is to study (singularities of) the quotient $\frac{p}{q}$ near $z_0$. What is the order of the poles, if any? Let's discuss a few instructive examples. 
\begin{enumerate}
    \item if $b_0 \neq 0$ so $q(z_0) \neq 0$. In particular $\frac{1}{q}$ is holomorphic near $z_0$ and we have 
    \begin{align*}
        \lim_{z-z_0}\frac{p(z)}{q(z)} = \frac{p(z_0)}{q(z_0)} = \frac{a_0}{b_0} \text{ exists, so }z_0 \text{ is a regular point for} \frac{p}{q} 
    \end{align*}
    \item if $a_0 = b_0 = 0$ but $b_1 \neq 0$ then 
    \begin{align*}
        \lim_{z-z_0}\frac{p(z)}{q(z)} = \lim_{z-z_0} \frac{a_1(z-z_0) + a_2(z-z_0)^2 + \cdots }{b_1(z-z_0) + b_2(z-z_0)^2 + \cdots } = \lim_{z-z_0}\frac{a_1 + a_2(z-z_0)^2 + \cdots }{b_1 + b_2(z-z_0)^2 + \cdots } = \\
        \frac{a_1}{b_1} \text{ exists hence } z_0 \text{ is a regular point for } f \text{ (l'Hôpital's rule)}
    \end{align*}
    Similarly, if $a_0=b_0=0$ and $a_1=b_1=0$ yet $b_2 \neq 0$
\end{enumerate}
\textbf{Example : } 
\begin{align*}
    f(z) = \frac{z}{sin(z)} \text{ where }z_0=0. \text{ using power series expansion of sin}
\end{align*}
\begin{align*}
    f(z) = \frac{z}{z-\frac{z^3}{3!}+\frac{z^5}{5!}+\cdots} = \frac{1}{1-\frac{z^2}{3!}+\frac{z^4}{5!}+\cdots}  
\end{align*}
Hence $\lim_{z \rightarrow 0} f(z) = 1$ exist, so $0$ regular point. \\
In summary either expand in power series and take out as many common powers of $z-z_0$ from the nominator and denominator as possible and take the limit or use l'Hôpital's rules. 